\section{SDD загвар -
Nogoolin}\label{doc6-sdd-ux437ux430ux433ux432ux430ux440---nogoolin}

\textbf{M2 · Тэнгис · 2026-09-21 · v0.2 · Бөглөх загвар}

SRS нь системээс шаардах үр дүнг, SDD нь түүнийг хэрэгжүүлэх аргыг
тодорхойлно. Lecture 02-ын M2 заалтын дагуу энэ бүтэц бэлтгэв. Доорх
асуултуудыг дараагийн ажлаар бөглөнө; бүрэн архитектурын баримт биш.

\subsection{1. Зорилго ба хамрах
хүрээ}\label{doc6-ux437ux43eux440ux438ux43bux433ux43e-ux431ux430-ux445ux430ux43cux440ux430ux445-ux445ux4afux440ux44dux44d}

\emph{{[}Бөглөх: ямар SRS хувилбар, FR/NFR ID-г энэ дизайн хангах вэ?
Хэн унших вэ?{]}}

Эхлэл: Nogoolin; \href{../../requirements/scope-charter.md}{Scope
Charter}, \href{../../requirements/requirements.md}{таван шаардлага}.
Уншигч нь хөгжүүлэгч Тэнгис, хичээлийн багш.

\subsection{2. Системийн орчин ба
бүрэлдэхүүн}\label{doc6-ux441ux438ux441ux442ux435ux43cux438ux439ux43d-ux43eux440ux447ux438ux43d-ux431ux430-ux431ux4afux440ux44dux43bux434ux44dux445ux4afux4afux43d}

\emph{{[}Бөглөх: системийн хил, гаднын оролцогч, бүрэлдэхүүн бүрийн
үүрэг ба холбоо. Context diagram нэм.{]}}

Эх материал: \href{https://github.com/Tengis01/Nogoolin}{Nogoolin public
repository}, web/admin, API, өгөгдөл/auth, shared validation;
\href{../../phase-0/README.md}{Phase 0 index}. Эдгээрийг бэлэн
архитектурын нотолгоо гэж үзэхгүй, холбогдох кодтой тулгана.

\subsection{3. Архитектурын
шийдвэр}\label{doc6-ux430ux440ux445ux438ux442ux435ux43aux442ux443ux440ux44bux43d-ux448ux438ux439ux434ux432ux44dux440}

\emph{{[}Бөглөх: шийдвэр, хувилбарууд, сонгосон шалтгаан, сул тал,
хангах requirement ID.{]}}

Карт: \textbf{Decision ID:} \nolinkurl{[D-XX]}; \textbf{Requirement:}
\nolinkurl{[FR/NFR-XX]}; \textbf{Сонголт:} \nolinkurl{[…]};
\textbf{Үндэслэл/эрсдэл:} \nolinkurl{[…]}; \textbf{Source:}
\nolinkurl{[…]}.

Next.js, Fastify, Supabase зэрэг технологийн сонголтыг
\href{../../../agent-context/DECISIONS.md}{шийдвэрийн эх}-ээс нягталж
энд оруулна. ``Хэрэглэгч юу хийж чадах вэ?'' гэсэн шаардлагыг
технологийн нэрээр орлуулахгүй.

\subsection{4. Өгөгдөл ба
хамгаалалт}\label{doc6-ux4e9ux433ux4e9ux433ux434ux4e9ux43b-ux431ux430-ux445ux430ux43cux433ux430ux430ux43bux430ux43bux442}

\emph{{[}Бөглөх: entity/холбоо, ownership, эрх, өгөгдлийн lifecycle;
шаардлагын ID ба ER эх.{]}}

Эхлэх жишээ: NFR-01-ийн хувийн өгөгдөл тусгаарлах шаардлагад session
identity, ownership filter, RLS/RBAC ямар үүрэгтэйг тайлбарлана.

\subsection{5. Интерфэйс ба ажиллах
урсгал}\label{doc6-ux438ux43dux442ux435ux440ux444ux44dux439ux441-ux431ux430-ux430ux436ux438ux43bux43bux430ux445-ux443ux440ux441ux433ux430ux43b}

\emph{{[}Бөглөх: интерфэйсийн contract холбоос, бүрэлдэхүүн хоорондын
дараалал, хэвийн/алдааны урсгал.{]}}

FR-01-ийн асуулга илгээх урсгалаас эхэлнэ. API field/schema-г нэг эхэд
хадгалж, энд холбоно; \href{../../phase-0/06-api-spec.yaml}{OpenAPI}-г
кодтой тулгах шаардлагатай.

\subsection{6. Чанарын шаардлагыг хангах
арга}\label{doc6-ux447ux430ux43dux430ux440ux44bux43d-ux448ux430ux430ux440ux434ux43bux430ux433ux44bux433-ux445ux430ux43dux433ux430ux445-ux430ux440ux433ux430}

\emph{{[}Бөглөх: NFR ID → дизайны арга → шалгах арга → эрсдэл.{]}}

Жишээ чиглэл: NFR-02-д fallback, NFR-03-д тестийн хаалт. Арга төлөвлөсөн
эсэх, хэрэгжсэн эсэх, шалгасан эсэхийг ялгана.

\subsection{7. Байршуулалт ба
хамаарал}\label{doc6-ux431ux430ux439ux440ux448ux443ux443ux43bux430ux43bux442-ux431ux430-ux445ux430ux43cux430ux430ux440ux430ux43b}

\emph{{[}Бөглөх: local/test/production орчин, компонент байрших газар,
build/deploy хамаарал, тохиргооны нэр.{]}} Нууц түлхүүрийн утга
бичихгүй. \href{../../phase-0/09-deployment.md}{Deployment эх}-ийг
ашиглана.

\subsection{8. Нээлттэй асуудал ба
төлөвлөгөө}\label{doc6-ux43dux44dux44dux43bux442ux442ux44dux439-ux430ux441ux443ux443ux434ux430ux43b-ux431ux430-ux442ux4e9ux43bux4e9ux432ux43bux4e9ux433ux4e9ux4e9}

\begin{longtable}[]{@{}
  >{\raggedright\arraybackslash}p{(\linewidth - 4\tabcolsep) * \real{0.2600}}
  >{\raggedright\arraybackslash}p{(\linewidth - 4\tabcolsep) * \real{0.5500}}
  >{\raggedright\arraybackslash}p{(\linewidth - 4\tabcolsep) * \real{0.1900}}@{}}
\toprule\noalign{}
\begin{minipage}[b]{\linewidth}\raggedright
Хэсэг
\end{minipage} & \begin{minipage}[b]{\linewidth}\raggedright
Бөглөх зүйл
\end{minipage} & \begin{minipage}[b]{\linewidth}\raggedright
Status
\end{minipage} \\
\midrule\noalign{}
\endhead
\bottomrule\noalign{}
\endlastfoot
§1--4 & Архитектурын тайлбар, шийдвэр, өгөгдөл/эрхийн дизайн & planned
(W04) \\
§5 & API contract нийцүүлэлт & planned (W05) \\
§2/4/5 зураг & Context, ER, sequence диаграмын эх & planned (W08) \\
§6--7 & CI/build шалгалтын нотолгоо & planned (W09) \\
\end{longtable}

Эдгээр нь дараагийн ажлын төлөвлөгөө; M2-д дээрх бөглөх бүтэц бэлэн
болсон. Шийдвэр өөрчлөгдвөл SRS ID, traceability, холбогдох диаграмыг
хамтад нь шинэчилнэ.
